\documentclass[11pt,a4paper]{scrartcl}
%Analogo di KOMA Script
% della classe article
%\usepackage{classicthesis-preamble}
\usepackage[nochapters,pdfspacing]{classicthesis}

\usepackage[utf8]{inputenc}
\usepackage[T1]{fontenc}
\usepackage{geometry}
\usepackage{booktabs}
\usepackage{array}
\usepackage{parskip}
\usepackage{titlesec}
\usepackage{enumitem}
\usepackage{amsmath}
\usepackage{amssymb}
\usepackage{hyperref}
\usepackage{xcolor}

\geometry{margin=2.5cm}
\usepackage{setspace}
\onehalfspace




\definecolor{rossoSapienza}{RGB}{130,36,51}

\titleformat{\section}{\large\bfseries\color{rossoSapienza}}{}{0em}{}[\titlerule]
\titleformat{\subsection}{\normalsize\bfseries\color{rossoSapienza}}{}{0em}{}
\titleformat{\subsubsection}{\normalsize\itshape}{}{0em}{}

\hypersetup{colorlinks=true, linkcolor=rossoSapienza, urlcolor=rossoSapienza}

\begin{document}
	\pagestyle{empty}
	
	% Title block
	\begin{center}
		{\LARGE\bfseries\color{rossoSapienza} Continuum Mechanics}\\[0.4em]
		{\large MSc in Nanotechnology Engineering}\\[0.2em]
		{\normalsize 6 CFU \quad|\quad English \quad|\quad 60 hours}
	\end{center}
	
	\vspace{1em}
	\hrule
	\vspace{1.5em}
	
	% Course Description
	\section{Course Description}
	
	The course provides a rigorous introduction to the foundations of Continuum Mechanics, structured around the four pillars of any continuum field theory: kinematics, stress, balance laws, and constitutive equations. These are developed with full mathematical rigor, then connected to the atomistic description of matter through two bridging frameworks: the Cauchy--Born rule, which shows how constitutive equations emerge from interatomic potentials, and the Irving--Kirkwood--Noll and Hardy formulations, which show how continuum fields such as stress and density emerge from a discrete atomic description.
	
%	The course serves as a formal foundation for the subsequent course in Molecular Dynamics and Atomistic Simulations, and connects directly to Surface Engineering, Nanomaterials, and Nanobiotechnology.
	
	% Prerequisites
	\section{Prerequisites}
	
	\begin{enumerate}[leftmargin=1.5em]
		\item Linear algebra and multivariable calculus
		\item Basic mechanics of materials (undergraduate level)
		\item Elements of classical physics
	%	\item Familiarity with index notation is helpful but not required
	\end{enumerate}
	
	% Module 1
	\section{Module 1 --- Mathematical Toolkit \normalfont\normalsize\itshape (8 h)}
	
	A focused review to establish a common formal language across the mixed student background.
	
	\begin{enumerate}[leftmargin=1.5em]
	%	\item Index notation and Einstein summation convention
		\item Tensor algebra
		\item Tensor calculus
	\end{enumerate}
	
	% Module 2
	\section{Module 2 --- Kinematics of Deformable Bodies \normalfont\normalsize\itshape (12 h)}
	\begin{enumerate}
		\item \textit{General Theory}. Deformation, displacement, deformation gradient. Change in  length, volume, and area. Homogeneous deformations. Rigid deformations. Strain measures. Polar decomposition.  Motion velocity, velocity gradient, stretching, and spin.
		\item  \textit{Infinitesimal theory}. Infinitesimal strain tensor and rotation tensor. Geometric changes.
	\end{enumerate}

	
	% Module 3
	\section{Module 3 --- Basic Mechanical Principless \normalfont\normalsize\itshape (12 h)}
	\begin{enumerate}
		\item \textit{Mass balance}
		\item  \textit{Force and moment balance}.  Internal actions: stress. Balance laws for linear and angular momentum. Ecternal power, kinetic energy, inertial power. Referential forms for the mechanical laws: Piola stress, expended power, second Piola stress. 
		\item \textit{Frame indifference}. Frame of reference. Frame-indifference principle. 
	%	\item \textit{Principle of Virtual Power}.  Force balance, stress, symmetry and frame-indifference.
	%	\item \textit{Energy balance}. The first law: balance of energy. 
	\end{enumerate}

	
	% Module 4
	\section{Module 4 --- Constitutive Theory \normalfont\normalsize\itshape (8 h)}

\begin{enumerate}
	\item \textit{Constitutive theories}.   Basic hypotheses for developing physically meaningful constitutive theories.
	\item \textit{Elastic Materials}. Consequence of frame-indifference, infinitesimal theory.
\end{enumerate}	

%		 Strain energy density function.
%		 Objectivity and material frame indifference.
%		 Neo-Hookean model as a worked example.
%		 Linear elasticity.

	
% Module 5
\section{Module 5 --- Atomistic-to-Continuum Bridging \normalfont\normalsize\itshape (17 h)}


\begin{enumerate}
	
	\item \textit{Failure of Classical Continuum Mechanics at the Nanoscale} (1 h).
	The continuum hypothesis and its length-scale assumptions. Surface-to-volume ratio
	scaling and its mechanical consequences. Size effects in elastic moduli.
	The Representative Volume Element (RVE) and its breakdown at the nanoscale.
	
	\item \textit{The Cauchy--Born Rule} (2 h).
	Bravais lattices and crystal kinematics: unit cell, primitive vectors, lattice
	deformation. The Cauchy--Born hypothesis: affine deformation of the lattice.
	From interatomic potential to strain energy density. Linearization: recovery of elastic constants. Limits of validity.
	
	\item \textit{Statistical Mechanics Primer} (3 h).
	Microstates and macrostates; phase space. Statistical ensembles: microcanonical,
	canonical, grand canonical. Time averages vs.\ ensemble averages; ergodic hypothesis.
	Liouville equation and conservation of phase space volume.
	
	\item \textit{Continuum Fields from Discrete Data} (11 h).
	
	\begin{enumerate}
		
		\item \textit{Irving--Kirkwood-Noll procedure} ($\sim$6 h).
		The fundamental problem: defining density, momentum, stress, and energy flux
		for a discrete set of atoms. Ensemble-averaged conservation equations derived
		from the Liouville equation. Stress tensor from pairwise forces: kinetic
		contribution and virial contribution.
		Non-uniqueness of the pointwise stress field; limitations of ensemble averaging.
		
		\item \textit{Hardy procedure} ($\sim$5 h).
		Localization function: from Dirac delta to smooth kernel.
		Pointwise continuum fields: density, momentum density,
		Hardy stress. Bond function and the geometric interpretation of stress
		localization. 
		
	\end{enumerate}
	
\end{enumerate}
	
	% Hour Distribution
%	\section{Hour Distribution}
%	
%	\begin{center}
%		\begin{tabular}{clc}
%			\toprule
%			\textbf{Module} & \textbf{Title} & \textbf{Hours} \\
%			\midrule
%			1 & Mathematical Toolkit            & 8  \\
%			2 & Kinematics of Deformable Bodies & 12 \\
%			3 & Stress and Balance Laws         & 12 \\
%			4 & Constitutive Theory             & 8  \\
%			5 & Atomistic-to-Continuum Bridging & 20 \\
%			\midrule
%			& \textbf{Total}                  & \textbf{60} \\
%			\bottomrule
%		\end{tabular}
%	\end{center}
%	

	
%	% Assessment
%	\section{Assessment}
%	
%	\begin{center}
%		\begin{tabular}{lcp{7cm}}
%			\toprule
%			\textbf{Component} & \textbf{Weight} & \textbf{Notes} \\
%			\midrule
%			Written exam       & 60\% & Tensor derivations and analytical problems \\
%			Take-home project  & 40\% & Analytical bridging problem --- e.g.\ Hardy stress
%			computation on an assigned atomic configuration,
%			or derivation of $W(\mathbf{F})$ via Cauchy--Born
%			for a given potential \\
%			\bottomrule
%		\end{tabular}
%	\end{center}
	
	% References
	\section{References}
	
Lecture Notes. Additional references:	
	
	\begin{enumerate}
		\item M.E. Gurtin, \textit{An Introduction to Continuum Mechanics}, Academic Press, 1982.
		\item M.E. Gurtin, E. Fried, L. Anand, \textit{The Mechanics and Thermodynamics of Continua}, Cambridge University Press, 2010.
		\item E.B. Tadmor, R.E. Miller, 	\textit{Modeling Materials}, Cambridge, 2011.
		\item C. Truesdell, \textit{A First Course in Rational Continuum Mechanics}, Academic Press, 1977.
	
	\end{enumerate}
	
	
\end{document}